\scanpage{31}
\mspage{29}
% archive sheet number: 30
% hand: Grothendieck
% sheet_kind: text
% element: text | upper, lines 1-6 | Proposition 6.2. Soit à nouveau pi
\uline{Proposition} 6.2. Soit \deleted{$\pi=\operatorname{Gal}(\overline K/K)$}, à nouveau $\pi$ comme dans $0.1$, \deleted{où le corps $k$} supposons que le corps résiduel $k$ soit tel que l'\inserted{\unsure{sous-}}extension \deleted{de} $\overline k$ engendrée par les racines $\ell^\nu$ ièmes de l'unité ($\nu\in\mathbb N$) soit infinie. Alors \deleted{\unsure{toute extension}} représentation $\ell$-adique de $\pi$ induit sur $I$ une repr. quasi-unipotente.

% element: text | middle, lines 7-13 | En effet, quitte à passer à une extension finie
En effet, quitte à passer à une extension finie $K'$ de $K$, et à remplacer $U$ par un ouvert, on peut supposer que l'action de $I$ se fait à travers $I(\ell)=I_t(\ell)=\mathbb Z_\ell(1)$, \deleted{l'hyp.} de sorte qu'on a en fait une représentation d'un quotient $\pi'$ de $\pi$, extension de $\pi_0$ par $\mathbb Z_\ell(1)\simeq\mathbb Z_\ell$. L'hypothèse faite sur $k$ signifie que l'hyp. de $6.1$ est vérifiée.

% element: text | lower, lines 14-15 | Corollaire 6.3 Toute représentation
\uline{Corollaire} 6.3 Toute représentation $\ell$-adique \deleted{\unsure{commutative}} \inserted{\unsure{ess.\textsuperscript{t}|virt.} abélienne} de $\pi$ induit une repr. \unsure{essentiellt|essentiellemt|essentiellement} triviale de $I$.

% element: text | lower, lines 16-17 | Cf II 1.12
Cf $\mathrm{II}\ 1.12$, qu'il y aurait lieu de reporter ici.
